% ===========================================================================
%  CAPITULO 5 — Editorial, Formação e Implantação (E01–E03)
% ===========================================================================
\chapter{Editorial, Formação e Implantação}
\label{cap:editorial}

As três perguntas deste bloco dependem, em grande parte, de material que não
está no repositório da plataforma. O que segue separa o que a plataforma
comprova do que o time editorial e o de implantação precisam anexar.

% ---------------------------------------------------------------------------
\section{E01 · Coleção didática e missões, por ano e bimestre}
\label{sec:e01}

\begin{resposta}{E01}{Matriz curricular integrada}{\foradorepo}
\begin{campos}
  \item[Plataforma] Cada missão é identificada por ano/série e unidade (1 a
    4), na grade de 12 séries $\times$ 4 unidades = 48 posições. Hoje há
    \textbf{14 missões publicadas}, cobrindo as 12 séries na unidade 1, mais
    a unidade 2 do 6º ano e a unidade 4 do 8º ano (Apêndice~\ref{ap:catalogo},
    com as habilidades BNCC de cada régua). Cada missão tem guia de aula com
    objetivos e habilidades BNCC (P01). Novas missões e versões são
    publicadas sem interromper o uso: cada versão tem endereço próprio,
    pode ser testada numa escola antes da rede e revertida (T09); toda
    missão nasce com interruptor de disponibilidade por rede
    (\adrr{ADR-87}).
  \item[Limitações] A plataforma não tem o vínculo com capítulos ou aulas do
    livro (só texto livre em alguns guias) nem a equivalência formal unidade
    = bimestre. A estrutura ``sete aulas com o livro e a oitava com a
    missão'' não aparece no repositório. Três missões não têm régua de
    avaliação por habilidade (Foto Conta, Rastro do Boato, Manifesto da IA).
  \item[Lacuna] Editorial: a matriz ano $\times$ bimestre com tema, objetivos,
    habilidades, capítulos/aulas e missão; a lista de componentes entregues
    (livro do aluno, do professor, recursos digitais); o processo de
    atualização curricular; amostras.
\end{campos}
\end{resposta}

% ---------------------------------------------------------------------------
\section{E02 · Orientações e materiais de apoio ao professor}
\label{sec:e02}

\begin{resposta}{E02}{Onde está cada orientação}{\parcial}
\begin{campos}
  \item[Na plataforma] (1) \textbf{Ficha da missão}: como funciona, preparo,
    competência e formato. (2) \textbf{Guia da aula} de toda missão
    publicada, com versão para impressão: objetivos, habilidades, preparo
    técnico e pedagógico, roteiro minuto a minuto, debate, mediação,
    síntese, rubrica, desdobramentos, cuidados, adaptações e dúvidas (P01).
    (3) \textbf{Exemplos de personalização}: o banco padrão de cada missão.
    (4) \textbf{Interpretação dos resultados}: critérios por nível no painel
    de habilidades. (5) \textbf{Ajuda} com assistente por voz e avatar para
    professores e gestores, cuja base de conhecimento é gerada do Mapa de
    Produto, dos guias e dos manifestos (\adrr{ADR-132}, \adrr{ADR-158}).
    (6) \textbf{Agente de IA do professor} sobre BNCC e o acervo de livros
    (parcial, P08).
  \item[Cobertura] O guia é obrigatório em toda missão publicada (um teste
    recusa publicar sem ele), então cobre as 14 missões.
  \item[Lacuna] Editorial: listar o que está no livro do professor e o que é
    entregue à parte, para o TR não exigir na plataforma o que está só no
    impresso, e vice-versa.
\end{campos}
\end{resposta}

% ---------------------------------------------------------------------------
\section{E03 · Implantação, formação e apoio às redes}
\label{sec:e03}

\begin{resposta}{E03}{Escopo de serviço}{\foradorepo}
\begin{campos}
  \item[Plataforma] Ferramentas que a implantação usa: uma infraestrutura
    isolada por contrato; carga de escolas, turmas e pessoas por OneRoster ou
    planilha, com relatório de erros (T03); criação de escolas, regionais e
    gestores pela cadeia de papéis (T04); interruptor de produto e missão por
    rede, regional ou escola (\adrr{ADR-87}); relatório da diretoria (T10); a
    Ajuda in-app.
  \item[Limitações] Não há no repositório plano de implantação, modelo de
    formação (público, carga horária, modalidade, certificação), canais de
    suporte técnico e pedagógico, regra de escalonamento nem entregáveis de
    fiscalização. Também não há registro de participação em formação na
    plataforma.
  \item[Lacuna] Implantação/Comercial: etapas e marcos de aceite, pré-requisitos
    da Secretaria, modelo de formação, canais e horários de suporte (ver
    T09), relatórios de fiscalização.
\end{campos}
\end{resposta}

\begin{paraotr}
A contratada deve apresentar, no início do contrato, plano de implantação com
etapas, responsáveis e marcos de aceite (carga cadastral validada, acesso de
gestores, formação inicial concluída, primeira missão aplicada por escola), e
entregar relatórios mensais de adoção (V1-01), de formação (participantes e
carga horária) e de atendimento (chamados por severidade e prazo de solução).
\end{paraotr}
